\section{AES encryption and IGE chaining}
\label{sec:aes}

The example in \texttt{example/aes\_ige/} adds a concrete block cipher to the
data-structure developments. It implements AES-128, AES-192, and AES-256,
including key expansion, block encryption and decryption, and
OpenSSL-compatible Infinite Garble Extension (IGE). Its proof coverage has
two parts: checked algebra for the concrete pure cipher definitions, and Iris
refinement contracts for individual executable kernels. Refinement of the
complete key-expansion routine, AES routines, and IGE loops remains outstanding.
Consequently, this section does not claim an end-to-end weakest-precondition
theorem for the extracted encryption API.

\subsection{Bytes, blocks, and the new primitives}

A byte is an integer satisfying \(0\leq x<256\). Four bytes form a column,
and four columns form a 16-byte block. The memory representation stores each
byte in one \texttt{Int64} word, in column order. This matches the existing
word-addressed raw heap; no packed-byte memory model is introduced.

Integer XOR combines bytes and round keys. Integer AND provides masking and
the high-bit test used by multiplication by two in the AES byte field:
\[
 d=(x+x)\mathbin{\mathrm{AND}}255,\qquad
 \mathsf{xtime}(x)=
 \begin{cases}
 d & \text{if }x\mathbin{\mathrm{AND}}128=0,\\
 d\mathbin{\mathrm{XOR}}27 & \text{otherwise}.
 \end{cases}
\]
For example, \(\mathsf{xtime}(128)=27\). Masking removes the ninth bit and
the conditional XOR implements reduction by the AES field polynomial.
The correspondence lemma \texttt{xtime\_correspondence} checks this arithmetic
implementation against the pure byte function for every byte.
\texttt{xtime\_wp} then proves the CJR call returns that value. Addition,
\texttt{BitAndOp}, and \texttt{BitXorOp} suffice; no multiplication, division,
or shift primitive is needed. The S-box and inverse S-box use concrete tables
lowered to branches, whose byte-domain correspondence is also checked.

\subsection{Concrete AES and inverse laws}

The files \texttt{bytes.v} and \texttt{block.v} define byte operations,
SubBytes and its inverse, ShiftRows and its inverse, MixColumns and its
inverse, and block XOR. The file \texttt{aes.v} defines the key schedule
and round composition. Keys of 16, 24, and 32 bytes select 10, 12, and 14
rounds respectively. Encryption begins with AddRoundKey; each intermediate
round applies SubBytes, ShiftRows, MixColumns, and AddRoundKey. The final
round omits MixColumns. Decryption reverses these transformations using the
same forward key schedule in reverse order.

Let \(E_k\) and \(D_k\) denote the concrete pure functions
\texttt{encrypt key} and \texttt{decrypt key}. On byte-valued blocks,
the checked theorems \texttt{decrypt\_encrypt} and
\texttt{encrypt\_decrypt} establish
\[
 D_k(E_k(b))=b,\qquad E_k(D_k(b))=b.
\]
The proofs compose the transformation inverse laws through the round list;
they do not assume an abstract cipher inverse axiom. Byte-range preservation
and injectivity/surjectivity are proved as well. The pure definitions are
total even for other key lists; the API supports only the three valid key
lengths. Inverse laws alone do not establish cryptographic security or
agreement with an external standard. Independent known-answer tests exercise
that agreement on the compiled implementation.

\subsection{IGE state, round trips, and chunking}

The file \texttt{ige.v} models the IV as a pair \((C_0,P_0)\): the previous
ciphertext block followed by the previous plaintext block. For each block,
\[
 C_i=E_k(P_i\mathbin{\mathrm{XOR}}C_{i-1})
        \mathbin{\mathrm{XOR}}P_{i-1},\qquad
 P_i=D_k(C_i\mathbin{\mathrm{XOR}}P_{i-1})
        \mathbin{\mathrm{XOR}}C_{i-1}.
\]
Both directions update the state to \((C_i,P_i)\). Thus decryption of an
encrypted message starts from the original IV, while continuation of the
same message starts from the updated IV.

Write \(\mathcal E_k(xs,v)\) and \(\mathcal D_k(xs,v)\) for the pure
IGE functions, each returning a block list and a final IV. For byte-valued
blocks and IV, \texttt{ige\_decrypt\_encrypt} states
\[
 \mathcal E_k(xs,v)=(cs,w)
 \quad\Longrightarrow\quad
 \mathcal D_k(cs,v)=(xs,w).
\]
\texttt{ige\_encrypt\_decrypt} proves the converse composition, including
the same final-IV agreement. Induction over the message uses the AES inverse
laws and cancellation of XOR to recover the block and the next state.

Length and byte-range preservation, injectivity, and surjectivity are checked.
Empty input returns the empty list and the original IV. The chunk theorem
\texttt{ige\_encrypt\_chunks} states that encrypting \(xs\mathbin{+\!+}ys\)
equals encrypting \(xs\), then encrypting \(ys\) with the first call's final
IV, concatenating the outputs and retaining the second call's IV.
\texttt{ige\_decrypt\_chunks} gives the corresponding decryption law.

\subsection{Executable contracts and caller obligations}

The CJR program is in \texttt{program.v}. The file \texttt{refinement.v}
proves S-box, inverse S-box, and byte-doubling contracts.
The file \texttt{memory.v} defines \texttt{block\_at}, a separating
conjunction owning the sixteen raw words of a block, and proves contracts
for block copy, ShiftRows/inverse ShiftRows, SubBytes/inverse SubBytes,
MixColumns, and block XOR. These contracts preserve source ownership and
replace the destination contents with the corresponding pure transformation.
For instance, \texttt{copy\_wp} has the shape
\[
 \mathsf{block\_at}(p,b)\sep\mathsf{block\_at}(q,old)
 \;\vdash\;
 \Wpre\,\mathsf{copy}(p,q)\,
 \{\lambda v.\;\pure{v=\Unit}\sep
   \mathsf{block\_at}(p,b)\sep\mathsf{block\_at}(q,b)\}.
\]
Separating ownership supplies the disjoint buffers needed by this contract.
These kernel results do not yet compose into a proof of the complete
executable encryption, decryption, or key-expansion API.

The extracted API begins with \texttt{aesExpandKey}. Key expansion reserves
240 words and writes
176/208/240 words for the three variants, returning 10/12/14 rounds;
invalid key length returns zero without changing the schedule.
The block routines, \texttt{aesEncryptBlock} and \texttt{aesDecryptBlock},
process sixteen words. The IGE routines, \texttt{aesIgeEncrypt} and
\texttt{aesIgeDecrypt}, process a nonnegative multiple of
sixteen logical bytes and updates a 32-word IV. All routines use a 64-word
workspace. Input, output, schedule, IV, and workspace buffers must be
disjoint, large enough, and supplied with byte-valued inputs and a matching
schedule and round count. These are caller obligations, not dynamic pointer
or capacity checks. Invalid round counts and invalid IGE message lengths
return false without changing output or IV. Empty messages preserve the IV
but still require a valid round count.

\subsection{Runtime evidence and limits}

The recorded runtime report contains fifteen passing cases: primitive
correspondence, exhaustive S-box inverse checks, NIST block known-answer
vectors for all three key sizes, OpenSSL IGE fixtures and published vectors,
chunk composition, empty input, invalid arguments, and buffer guards.
The executable cipher never calls OpenSSL; the reference fixtures are
generated independently and record library provenance.

The example also benchmarks extracted Cangjie against OpenSSL's native C
\texttt{AES\_ige\_encrypt} API. Complete outputs and final IVs are checked
before timing. Both harnesses reuse expanded schedules, reset IVs, and
checksum outputs inside the timed region, with three processes and seven
samples per workload. Throughput is measured per logical byte; forward
key-expansion timing includes its checksum. OpenSSL inverse-schedule setup
is excluded from decryption throughput and is not separately measured.
The offline report at \texttt{example/aes\_ige/docs/report.html} links the
measurements, comparison charts, fixtures, and proof declarations.

The comparison combines word-per-byte versus packed-byte representations,
cipher algorithms, library optimization, and compiler effects. It does not
isolate language overhead or establish use of hardware AES instructions.
The implementation provides no padding, authentication, or in-place
operation. Its branch-based S-box is not constant-time, and no side-channel
or cryptographic security theorem is claimed. Kernel proof checking
validates the stated algebra and refinement contracts; full executable
refinement and the printer/compiler correspondence remain separate limits.
